\chapter{Introduction}

\section{eBPF Background}
Extended Berkeley Packet Filter (eBPF) allows user-space programs to safely execute custom logic inside the Linux kernel without needing custom modules. Because these programs run in ring-0, a single vulnerability could easily crash the system or lead to privilege escalation. To prevent this, every eBPF program must pass the \textbf{BPF verifier}, a static analysis engine inside the kernel that formally checks the program for memory safety and type correctness before JIT compilation. 

This report documents an analysis of four significant logic bugs affecting the verifier, comparing their root causes, exploitation primitives, and the protective mechanisms that mitigate them.

\section{Research Objectives}
This project is part of the Security Verification and Testing course. The main goal is to analyze the patches, understand the vulnerabilities at a deep technical level, and develop proof-of-concept code for several verifier CVEs. While there are many vulnerabilities available for study, I have selected the following ones for this research:
\begin{itemize}
    \item \textbf{CVE-2021-3600}: A 32-bit truncation bug during the verifier's div/mod instruction fixup, leading to bounds tracking confusion and Out-Of-Bounds (OOB) memory access.
    \item \textbf{CVE-2021-3444}: The sister bug of CVE-2021-3600, targeting the mod32 destination register truncation. Exploited on kernel 4.19.19 (where \texttt{alu\_limit} is absent) to achieve a \textbf{full unprivileged local privilege escalation} via ret2usr.
    \item \textbf{CVE-2023-52452}: An integer tracking flaw regarding variable-offset stack writes, which incorrectly calculates the depth of stack frames and causes stack corruption.
    \item \textbf{CVE-2024-26589}: A missing bounds check inside the restricted pointer arithmetic mechanisms which allows arbitrary offset-addition to the \texttt{flow\_keys} structures, leading to stack OOB.
\end{itemize}

\section{Research Methodology}
The exploitation process is organized into distinct, verifiable \textit{Milestones}, starting from initial vulnerability confirmation (such as information leaks) up to more advanced exploitation stages and analyzing the technical feasibility of different attack vectors. 

This modular structure ensures that the research methodology remains both robust and extensible for future vulnerability analyses, allowing for the seamless integration of additional CVE studies regardless of the specific subsystem or bug class involved.

\section{Common Experimental Framework}

To improve reproducibility and keep the methodology scalable as new CVEs are added, experiments are conducted using a shared framework and CVE-specific parameterization. This section defines what is common across analyses; concrete values (kernel release, config variants, and mitigation profile) are documented inside each CVE chapter.

\begin{table}[H]
\centering
\small
\begin{tabular}{@{}ll@{}}
\toprule
\textbf{Component} & \textbf{Common baseline} \\
\midrule
Build workflow & Buildroot-based pipeline orchestrated by project scripts \\
Virtualization & QEMU/KVM virtual testbed \\
Root filesystem strategy & Minimal Buildroot image with CVE overlays \\
Kernel strategy & Per-CVE vulnerable target kernel selected for validation \\
eBPF core requirements & \texttt{CONFIG\_BPF\_SYSCALL=y}, \texttt{CONFIG\_BPF\_JIT=y} \\
Debug and observability & Kernel symbols/debug options enabled when needed by the CVE \\
Mitigation profile handling & Controlled and documented per scenario (privileged/unprivileged) \\
Compiler/toolchain strategy & Static user-space binaries for deterministic VM execution \\
Execution model & Milestone-driven validation and incremental exploit development \\
\bottomrule
\end{tabular}
\caption{Shared experimental framework used across CVE analyses}
\end{table}

This separation between common framework and per-CVE specifics keeps the report coherent, extensible, and directly reusable for future verifier vulnerability studies.
